\documentclass[11pt]{article}
\usepackage[a4paper,margin=2.2cm]{geometry}
\usepackage{amsmath,amssymb,booktabs,array}
\usepackage[T1]{fontenc}
\usepackage[utf8]{inputenc}

\title{Exact Symbolic Differences Between Watson--Aliev (1985) and Aliev (1986) for Linear-Molecule Quartic State Dependence}
\author{}
\date{2026-03-28}

\begin{document}
\maketitle

\section*{Statement}

This note records the exact symbolic conclusions obtained from a direct reduction of the printed Watson--Aliev formulas to the linear-molecule branch, and compares them with the compact formulas printed by Aliev for linear molecules.

\section*{Definitions}

Let
\[
U_{nn'}^{\mathrm{A86}},\qquad V_{nn'}^{\mathrm{A86}}
\]
denote the compact linear descendants printed in Aliev (1986), and let
\[
H_{24}^{\mathrm{bare}}(UV)
\]
denote the explicit bare \(U/V\) sector of Table~V in Watson--Aliev (1985).

The linear observable is written as
\[
D_v = D_J - \sum_n \beta_n\left(v_n+\frac12\right)-\sum_t \beta_t(v_t+1).
\]

\section*{Exact Results}

\subsection*{1. Bare \(H_{24}(UV)\) gives only \(V\)}

After exact specialization of
\[
M_{nt}=\frac{X_{nt}-X^{nt}}{4},\qquad
N_{nt}=\frac{X_{nt}+X^{nt}}{4},
\]
the bare Table-V block
\[
-8\Big[(2\omega_t+\omega_n+\omega_{n'})M_{nt}M_{n't}
+
(2\omega_t-\omega_n-\omega_{n'})N_{nt}N_{n't}\Big]
\]
reduces identically to
\[
-4(\omega_{n'}-\omega_n)\,V_{nn'}^{(0)},
\]
where \(V_{nn'}^{(0)}\) is the no-\(r_{nn'}\) part of the compact Aliev \(V\) channel.

Therefore
\[
\boxed{H_{24}^{\mathrm{bare}}(UV)\Rightarrow V\text{-only}.}
\]

\subsection*{2. The compact \(U\) channel is not produced by the printed bare kernel}

The compact linear combination contains
\[
(\omega_n+\omega_{n'})U_{nn'}^{\mathrm{A86}}
+
(\omega_{n'}-\omega_n)V_{nn'}^{\mathrm{A86}},
\]
but the explicit bare \(H_{24}(UV)\) sector supplies only the second channel.

Thus
\[
\boxed{U_{nn'}^{\mathrm{A86}}\not\subset H_{24}^{\mathrm{bare}}(UV).}
\]

\subsection*{3. The missing \(U\) channel cannot come from the printed \(X/F\) or \(R\) pieces of bare \(H_{24}\)}

Using the exact grading in \((C,\zeta,k_4)\),
\[
\Delta := U\text{-channel residue} \sim (0,2,0).
\]

For the printed \(X/F\) sector one finds only
\[
(4,2,2,0)\quad\text{and}\quad(2,1,1,1)
\]
in the corresponding \((B,\zeta_1,\zeta_2,r)\) bookkeeping, so that sector cannot generate a pure \(B^2\zeta^2\) term.

For the printed \(R\)-type terms of bare \(H_{24}\), the possible grades are
\[
(4,0,0),\ (2,0,1),\ (2,2,0),\ (2,0,0),\ (1,0,1),\ (1,2,0),\ (3,1,0),\ (2,1,0),
\]
and none equals
\[
(0,2,0).
\]

Therefore
\[
\boxed{\text{the printed bare }X/F\text{ and }R\text{ terms do not generate }U.}
\]

\subsection*{4. The only constructive source for \(U\) is \(i[S_{11}^{(R)},H_{04}]\)}

From Eq.~(106),
\[
\widetilde H_{24}-H_{24}
=
i[S_{03},H_{12}]
+i[S_{11}^{(R)},H_{04}]
+i[S_{13}^{(R)},H_{02}].
\]

Degree analysis shows:
\[
i[S_{03},H_{12}] \not\sim (0,2,0),
\]
\[
i[S_{11}^{(R)},H_{04}] \sim (0,2,0)\ \text{or}\ (2,0,0),
\]
while \(i[S_{13}^{(R)},H_{02}]\) depends on the already unresolved off-diagonal \(H_{24}\) sector.

Hence
\[
\boxed{\text{the only constructive source of }U\text{ is }i[S_{11}^{(R)},H_{04}].}
\]

\subsection*{5. The regular branch of \(S_{11}^{(R)}\) kills \(U\)}

For the accidental \(xy\) block,
\[
s_{xy}^{kl}=\frac{X_{xy}^{kl}}{4(B_y-B_x)}.
\]

In the exact linear limit the transverse quartic block must be cylindrically symmetric:
\[
H_{04}^{(xy)}=\lambda(J_x^2+J_y^2)^2,
\]
so that
\[
\tau_{xxxx}=\tau_{yyyy}=4\lambda,\qquad \tau_{xxyy}=\frac{4\lambda}{3},
\]
and therefore
\[
-\tau_{xxxx}+3\tau_{xxyy}=0,\qquad \tau_{yyyy}-3\tau_{xxyy}=0.
\]

Thus
\[
i[J_z,H_{04}^{(xy)}]=0
\]
in the exact linear limit.

If one imposes the regular branch
\[
X_{xy}=4(B_y-B_x)\widehat U_{xy},
\]
then
\[
\lim_{B_x\to B_y} i[S_{11}^{(R)},H_{04}] = 0.
\]

Therefore
\[
\boxed{\text{the regular branch of }S_{11}^{(R)}\text{ annihilates }U.}
\]

\subsection*{6. A singular branch is necessary}

Let
\[
\Delta:=B_y-B_x.
\]
Writing
\[
\tau_{xxxx}=3\sigma+\Delta a,\qquad
\tau_{yyyy}=3\sigma+\Delta b,\qquad
\tau_{xxyy}=\sigma+\Delta c,
\]
gives
\[
-\tau_{xxxx}+3\tau_{xxyy}=\Delta(-a+3c),
\]
\[
\tau_{yyyy}-3\tau_{xxyy}=\Delta(b-3c).
\]

With
\[
s_{xy}=\frac{X_{xy}}{4\Delta},
\]
the limit is finite:
\[
\lim_{\Delta\to 0}s_{xy}(-\tau_{xxxx}+3\tau_{xxyy})
=\frac{X_{xy}}{4}(-a+3c),
\]
\[
\lim_{\Delta\to 0}s_{xy}(\tau_{yyyy}-3\tau_{xxyy})
=\frac{X_{xy}}{4}(b-3c).
\]

Hence
\[
\boxed{U\neq 0\ \text{requires the singular branch}\ S_{11}^{(R)}\sim (B_y-B_x)^{-1}.}
\]

\subsection*{7. Eq.~(101) is complete as printed, and its \(\zeta^2\) sector still does not match Aliev's compact \(U\)}

From the vector PDF, Eq.~(101) contains exactly:
\[
\frac{3}{8}\omega_k\omega_l \sum_\gamma
\frac{C_k^{\alpha\gamma}C_l^{\beta\gamma}+C_l^{\alpha\gamma}C_k^{\beta\gamma}}{B_\gamma^e}
+\frac12\sum_m k'_{klm}C_m^{\alpha\beta}
\]
plus exactly two \(\zeta^2\) lines:
\[
K_1(\zeta_{mk}^{\alpha}\zeta_{ml}^{\beta}+\zeta_{ml}^{\alpha}\zeta_{mk}^{\beta}),
\]
\[
K_2(\zeta_{mk}^{\alpha}\zeta_{ml}^{\beta}+\zeta_{ml}^{\alpha}\zeta_{mk}^{\beta}),
\]
with
\[
K_1=
-\frac14 B_\alpha^e B_\beta^e
\frac{(\omega_m-\omega_k)(\omega_m-\omega_l)(2\omega_m+\omega_k+\omega_l)}
{\omega_m(\omega_m+\omega_k)(\omega_m+\omega_l)(\omega_k\omega_l)^{1/2}},
\]
\[
K_2=
-\frac14 B_\alpha^e B_\beta^e
\frac{(\omega_m+\omega_k)(\omega_m+\omega_l)(2\omega_m-\omega_k-\omega_l)}
{\omega_m(\omega_m-\omega_k)(\omega_m-\omega_l)(\omega_k\omega_l)^{1/2}}.
\]

The crop shows the closing bracket and equation number, so no third \(\zeta^2\) term is omitted.

The \(\zeta^2\) kernel extracted from Eq.~(101) is therefore complete and equals
\[
K_{101}^{(\zeta^2)}=-\frac14(K_1+K_2),
\]
but
\[
K_{101}^{(\zeta^2)}\neq K_U^{\mathrm{A86}},
\]
where
\[
K_U^{\mathrm{A86}}=
\frac{(\omega_k+\omega_m)^2(\omega_l-\omega_m)^2+(\omega_k-\omega_m)^2(\omega_l+\omega_m)^2}
{8\omega_m\sqrt{\omega_k\omega_l}(\omega_k^2-\omega_m^2)(\omega_l^2-\omega_m^2)}.
\]

Therefore
\[
\boxed{\text{Eq.~(101) as printed does not reproduce the compact Aliev }U\text{ kernel.}}
\]

\subsection*{8. The discrepancy survives in the final observables}

From the explicit production code used to build the linear observables, the \(U/V\) entry point is

\[
\mathrm{Term7}_{\parallel}
=-16\Big[U_{ij}^2(\omega_i+\omega_j)+V_{ij}^2(\omega_j-\omega_i)\Big],
\]

\[
\mathrm{Term7}_{\perp}
=-16\Big[U_{tj}^2(\omega_t+\omega_j)+V_{tj}^2(\omega_j-\omega_t)\Big].
\]

Hence
\[
\frac{\partial \mathrm{Term7}_{\parallel}}{\partial U_{ij}}
=-32U_{ij}(\omega_i+\omega_j)\neq 0,
\]
\[
\frac{\partial \mathrm{Term7}_{\perp}}{\partial U_{tj}}
=-32U_{tj}(\omega_t+\omega_j)\neq 0.
\]

Thus the unresolved \(U\)-channel is not projected away in the final observables:
\[
\boxed{
\beta_{\parallel}\ \text{and}\ \beta_{\perp}\ \text{both depend essentially on }U.
}
\]

\section*{Final Comparison}

\begin{center}
\begin{tabular}{>{\raggedright\arraybackslash}p{3.2cm} >{\raggedright\arraybackslash}p{4.7cm} >{\raggedright\arraybackslash}p{4.7cm} >{\raggedright\arraybackslash}p{3.2cm}}
\toprule
Object & Watson--Aliev (1985), printed form & Aliev (1986), printed compact form & Exact relation \\
\midrule
\(V_{nn'}\) & Obtained from bare \(H_{24}(UV)\) & Present explicitly & Derivable \\
\(U_{nn'}\) & Not obtained from bare \(H_{24}\); requires \(i[S_{11}^{(R)},H_{04}]\) and the singular branch & Present explicitly & Not derivable from printed `Paper1` \\
Eq.~(101) \(\zeta^2\) kernel & Complete and explicit & Compact \(U\) kernel implicit in \(U_{nn'}\) & Noncoincident \\
\(\beta_{\parallel}\) & Depends on \(U^2+V^2\) & Depends on compact \(U,V\) & Not derivable from printed `Paper1` \\
\(\beta_{\perp}\) & Depends on \(U^2+V^2\) & Depends on compact \(U,V\) & Not derivable from printed `Paper1` \\
\bottomrule
\end{tabular}
\end{center}

\section*{Conclusion}

The exact symbolic reduction of the printed Watson--Aliev formulas to the linear case gives:

\[
\boxed{
\text{the compact Aliev formulas for }U_{nn'},\ \beta_{\parallel},\ \beta_{\perp}
\text{ are not derivable from the printed form of Watson--Aliev.}
}
\]

The discrepancy is exact, not numerical. It is localized in the \(U\)-channel and survives all the way to the final observables.

\end{document}
